%! Graphing Radical Functions \& Transformations — Unit 6, Lesson 6.4 — Warm-Up
\documentclass[10pt]{article}
\usepackage{algebra2-article}
\usepackage{algebra2-boxes}
\newcommand{\TallMath}[1]{$\displaystyle #1\rule[-1.4em]{0pt}{3.2em}$}

\begin{document}

\pageheader{Unit 6, Lesson 6.4}{Warm-Up}
\namedateperiod

\vspace{0.10in}

\begin{notesbox}{Getting Ready: Skills We'll Use Today}
\small Three skills you already have. Today they become one method. Work quickly.

\vspace{4pt}
\textbf{1. Shifting a graph (Units 2--4).} Describe how each graph moves away from its parent, and
give the point the graph is built around.
\begin{center}
\renewcommand{\arraystretch}{1.6}
\begin{tabularx}{\linewidth}{@{}l X X@{}}
 & \textbf{How it moves from the parent} & \textbf{Vertex} \\
\hline
(a) \; $y=(x-4)^{2}+1$ & \blank{4.6cm} & \blank{2.0cm} \\
(b) \; $y=|x+2|-3$ & \blank{4.6cm} & \blank{2.0cm} \\
\end{tabularx}
\end{center}
One of the two numbers in each equation moves the graph in the \emph{same} direction as its sign,
and the other moves it in the \emph{opposite} direction. Which is which? \writeline

\vspace{4pt}
\textbf{2. Yesterday's domains, one step further (Lesson 6.0).} Fill in the row.
\begin{center}
\renewcommand{\arraystretch}{1.6}
\begin{tabularx}{\linewidth}{@{}l X X X@{}}
\hline
\textbf{Function} & \textbf{Solve radicand} $\ge 0$ & \textbf{Domain} & \textbf{Endpoint} \\
\hline
$g(x)=\sqrt{x-3}$ & \blank{2.4cm} & \blank{2.4cm} & \blank{2.0cm} \\
$h(x)=\sqrt{x+5}$ & \blank{2.4cm} & \blank{2.4cm} & \blank{2.0cm} \\
\hline
\end{tabularx}
\end{center}
Compare the number you see \emph{inside} each radical with the first coordinate of that function's
endpoint. What do you notice? \writeline

\vspace{4pt}
\textbf{3. Pulling a number out (Lesson 6.2).} Simplify. Assume $x\ge 0$.
\begin{center}
\renewcommand{\arraystretch}{1.6}
\begin{tabular}{@{}l l l@{}}
(a) \; $\sqrt{9x}=$ \blank{1.8cm} & (b) \; $\sqrt{25x}=$ \blank{1.8cm} &
(c) \; $\sqrt[3]{8x}=$ \blank{1.8cm} \\
\end{tabular}
\end{center}
Look hard at (a). You just claimed that $\sqrt{9x}$ and $3\sqrt{x}$ are the same expression. But one
of them is the parent $\sqrt{x}$ with its \emph{input} multiplied by $9$, and the other is the parent
with its \emph{output} multiplied by $3$. Can both of those really be the same graph?
\blank{2.4cm}

\end{notesbox}

\end{document}
